\section{Marco Actor-Critic}

\subsection{Mejora del gradiente de política con la función de ventaja}

Se sustituye el reward-to-go por $A^\pi$:
$$\nabla_\theta J(\theta) \approx \frac{1}{N}\sum_{i=1}^{N}\sum_{t=1}^{T} \nabla_\theta \log \pi_\theta(a_{i,t} \mid s_{i,t}) \cdot \hat{A}^\pi(s_{i,t}, a_{i,t})$$

Una mejor estimación de $\hat{A}$ $\Rightarrow$ un gradiente con menos ruido.

\subsection{Cómo estimar la función de valor}

Basta con ajustar $\hat{V}^\pi$, porque:
$$\hat{A}^\pi(s_t, a_t) = r(s_t, a_t) + \gamma \hat{V}^\pi(s_{t+1}) - \hat{V}^\pi(s_t)$$

Aquí se usa el $s_{t+1}$ obtenido por muestreo para aproximar la esperanza.

\subsection{Métodos para ajustar la función de valor}

\textbf{Método 1: regresión de Monte Carlo}

$$\min_\phi \sum_{s_t \in \mathcal{D}} \left\| \hat{V}^\pi_\phi(s_t) - \sum_{t'=t}^{T} r(s_{t'}, a_{t'}) \right\|^2$$

Se usa directamente la recompensa acumulada muestreada como target. Sin sesgo, pero con alta varianza.

\textbf{Método 2: actualización Bootstrapped / TD}

$$\min_\phi \sum_{(s,a,s') \in \mathcal{D}} \left\| \hat{V}^\pi_\phi(s) - \left(r(s,a) + \gamma \hat{V}^\pi_{\phi'}(s')\right) \right\|^2$$

Se usa la propia estimación como target (bootstrap). Baja varianza, pero con sesgo.

\textbf{Método 3: retorno de n pasos}

Combina los dos métodos anteriores: usa las recompensas reales de $n$ pasos más bootstrap:
$$y_{i,t} = \sum_{t'=t}^{t+n-1} r(s_{i,t'}, a_{i,t'}) + \gamma^n \hat{V}^\pi(s_{i,t+n})$$

\begin{knowledgebox}{Bias-Variance Tradeoff}
\begin{itemize}
    \item Monte Carlo: sin sesgo, alta varianza.
    \item TD (bootstrap de 1 paso): con sesgo (porque $\hat{V}$ no es exacta), baja varianza.
    \item n-step: ofrece un equilibrio ajustable entre ambos.
\end{itemize}
\end{knowledgebox}

\subsection{Arquitectura y flujo de datos}

\begin{figure}[H]
\centering
\includegraphics[width=0.9\textwidth]{slide-03.jpg}
\caption{Diagrama del flujo de datos de Actor-Critic en las diapositivas de la clase: la política (Actor) produce las acciones, y el critic aporta la estimación de la ventaja para retroalimentar el gradiente.}
\end{figure}

En la figura se puede ver un ciclo típico de Actor-Critic: se generan acciones a partir de la política actual, que se guardan en el conjunto de experiencia; el critic usa los datos de experiencia para ajustar la función de valor o los valores Q; la estimación de Advantage se envía de regreso al Actor para la actualización del gradiente. Todo el proceso es asíncrono, pero mantener con precisión fina la consistencia de los datos (fixed target network, actualizaciones retrasadas) es clave para la estabilidad.

\subsection{El critic como baseline y término de regularización}

El critic no solo calcula la ventaja, sino que también funge como baseline y puede mejorar la robustez mediante términos de regularización. Por ejemplo, la regularización L2, el value clipping y el moving average del bootstrap target son prácticas comunes para evitar el critic drift.

\begin{warningbox}{El colapso del entrenamiento causado por el critic drift}
Si el critic se entrena demasiado rápido, arrastra la función de valor hacia una región de sobreestimación, lo que hace que el advantage sea siempre positivo; el Actor tenderá a preferir ciertas acciones y convergerá rápidamente a una política degradada. Controlar la tasa de aprendizaje del critic y agregar una target network son técnicas habituales para evitar este "pessimistic collapse".
\end{warningbox}

\begin{knowledgebox}{Probabilistic baseline vs learned critic}
El gradiente de política tradicional usa el episode average reward como baseline, mientras que Actor-Critic aprende directamente una función parametrizada, más ajustada al estado actual. La ventaja de un critic aprendido es que puede manejar dependencias de largo plazo (como el retorno de n pasos con bootstrap) y corregir activamente el sesgo, lo que mejora la eficiencia de muestreo.
\end{knowledgebox}

\subsection{Generalized Advantage Estimation}

GAE combina los retornos de múltiples pasos mediante un núcleo de decaimiento exponencial, con la siguiente fórmula:
$$\hat{A}^\text{GAE}_{t} = \sum_{l=0}^{\infty} (\gamma \lambda)^l \delta_{t+l}, \qquad \delta_t = r_t + \gamma V(s_{t+1}) - V(s_t)$$
donde $\lambda \in [0,1]$ controla el bias-variance tradeoff. Con $\lambda=0$ se reduce al TD de 1 paso, y con $\lambda=1$ se reduce a Monte Carlo.

\begin{importantbox}{La estrategia de programación de GAE}
En la práctica, hacer anneal de $\lambda$ de manera gradual, de un valor bajo hacia 1, permite mantener una varianza baja al inicio del entrenamiento y absorber poco a poco, conforme avanza la convergencia, la información de dependencias de largo plazo. Combinado con una target network y el decay de $\gamma$, esto también evita que el critic se sobreajuste al reward de corto plazo.
\end{importantbox}

\subsection{Flujo completo de Actor-Critic}

\begin{enumerate}
    \item Ejecutar la política para recolectar datos
    \item Ajustar la función de valor $\hat{V}^\pi$
    \item Estimar la ventaja $\hat{A}^\pi$
    \item Actualizar la política $\pi_\theta$ con el gradiente de política
    \item Repetir
\end{enumerate}

\subsection{Resumen de la sección}

Actor-Critic aprende una función de valor para estimar la ventaja, en sustitución del reward-to-go muestreado, muy ruidoso, del gradiente de política, con lo que reduce de forma considerable la varianza del gradiente.

